\section{Cleaning-based quantitative rounding}
\label{sec:quantitative}

We now turn from exact equivalence to approximate symmetry.  Estimating one
reduced operator would lose the exponentially small nonidentity signal of a
minimum-support marginal.  The global argument avoids that loss.  Choose one
AME party as a logical input.  The other \(2m-1\) parties form an
\([[2m-1,1,m]]_q\) stabilizer code, and the product of their local factors
approximately implements the transpose-inverse of the factor on the chosen
party.  Three correctable regions then test that logical action by a nested
commutator.  Indeed, distance \(m\) makes every set of at most \(m-1\)
erased physical parties correctable, and the \(2m-1\) physical parties split
into regions of sizes \(m-1,m-1,1\).

The exact mechanism is standard: the qubit stabilizer
cleaning lemma is due to Bravyi and Terhal
\cite[Lemma~1]{BravyiTerhal2009}; in the additive prime-power setting used
here, the same cleaning step is the elementary symplectic statement recorded
in the proof of Lemma~\ref{lem:quantitative-cleaning-commutator}.
Nested-commutator localization underlies the
Bravyi--K\"onig hierarchy theorem \cite{BravyiKonig2013}, and
Pastawski--Yoshida prove directly that a transversal logical unitary on three
cleanable regions is Clifford \cite[Lemma~5]{PastawskiYoshida2015}.
Disjointness gives a complementary exact-code hierarchy bound
\cite[Theorem~5]{JochymOConnorKubicaYoder2018}.  The point below is therefore
not a new qualitative transversal-gate obstruction.  It is the robust step:
an implementation with leakage error \(\varepsilon\) forces every nested
logical Weyl commutator within \(8\varepsilon\) of a scalar, after which a
Fourier argument gives an explicit Clifford distance.

For an encoding isometry \(V:\mathbb C^q\to\mathcal H\), write
\[
 E(T,L)=q^{-1/2}\lVert TV-VL\rVert_{\mathrm{HS}}.
\]
It satisfies
\(E(TS,LM)\leq E(T,L)+E(S,M)\) and
\(E(T^\dagger,L^\dagger)=E(T,L)\).

\begin{lemma}[Quantitative cleaning commutator]
\label{lem:quantitative-cleaning-commutator}
Let \(V\) encode one logical qudit in a stabilizer code whose physical
parties have a partition into three correctable regions.  Let \(T\) be
transversal, and suppose \(E(T,L)\leq\varepsilon\).  For all logical Weyl
operators \(P,Q\),
\[
 q^{-1/2}\min_{|z|=1}
 \lVert [[L,P],Q]-zI\rVert_{\mathrm{HS}}
 \leq8\varepsilon,                                  \tag{5.1}
\]
where \([A,B]=ABA^\dagger B^\dagger\).
\end{lemma}

\begin{proof}
Let \(R_1,R_2,R_3\) be the assumed correctable partition.  The stabilizer
cleaning lemma supplies a representative \(\widetilde P\) of \(P\) outside
\(R_1\) and a representative \(\widetilde Q\) of \(Q\) outside \(R_2\).
For reference, cleaning here is the elementary symplectic statement that
the restriction of a normalizer label to a correctable region lies in the
restriction of the stabilizer label space; subtracting that stabilizer
removes the restriction.

Transversality makes \([T,\widetilde P]\) supported outside \(R_1\), so
\(C=[[T,\widetilde P],\widetilde Q]\) is supported on
\(R_1^c\cap R_2^c=R_3\).  The composition rules give
\[
 E(C,[[L,P],Q])\leq4\varepsilon.                      \tag{5.2}
\]
Knill--Laflamme correction on \(R_3\) gives
\(V^\dagger C V=cI\).  Hence the logical nested commutator lies within
\(4\varepsilon\) of \(cI\).  Its normalized Hilbert--Schmidt norm is one,
so \(|c|\geq1-4\varepsilon\).  If \(c=0\), then \(\varepsilon\geq1/4\), and
(5.1) is immediate because its left-hand side is at most
\(2\leq8\varepsilon\).  Otherwise, replacing \(c\) by \(c/|c|\) proves (5.1).
Notice that \(C\) need not preserve the code.  Only its scalar compression
is used, while (5.2) controls leakage.
\end{proof}

The remaining step is finite-dimensional Fourier concentration.  It is
stated separately because it is the point where scalar commutators become
Clifford axes.

\begin{lemma}[Nested Weyl commutators round to a Clifford]
\label{lem:nested-weyl-rounding}
Let \(L\in U(q)\), \(q=p^e\), and suppose for all Weyl labels \(v,w\)
that
\[
 q^{-1/2}\min_{|z|=1}
 \lVert [[L,W_v],W_w]-zI\rVert_{\mathrm{HS}}\leq\eta.
\]
If
\[
 \eta<\tau_p:=\min\left\{\frac13,
       \frac{\sin(\pi/p)}{2\sqrt2}\right\},
\]
then some additive one-qudit Clifford \(K\) satisfies
\[
 q^{-1/2}\min_{|z|=1}\lVert L-zK\rVert_{\mathrm{HS}}\leq\eta.
\]
\end{lemma}

\begin{proof}
Put \(A_v=LW_vL^\dagger\).  Conjugation of \(W_w\) by \(W_v\) changes
only its phase, so every \([A_v,W_w]\) is within \(\eta\) of a scalar.
We use the following Fourier observation twice.  If
\(A=\sum_a c_aW_a\) is unitary and every \([A,W_w]\) is within \(\eta\)
of a scalar, then \(A\) lies within \(\eta\) of one Weyl axis.  Indeed,
with \(p_a=|c_a|^2\), the commutator estimate gives
\[
 \left|\sum_a p_a\chi_w(a)\right|\geq1-\eta^2/2.
\]
Character orthogonality, averaged over \(w\), yields
\(\sum_ap_a^2\geq(1-\eta^2/2)^2\).  Thus
\(\max_ap_a\geq(1-\eta^2/2)^2\), which is precisely the asserted
normalized distance bound.

Apply the observation to each \(A_v\), and call its unique nearby Weyl
label \(F(v)\).  The exact identity
\(A_{v+u}\in\mathbb C A_vA_u\) puts the axes labelled by \(F(v+u)\) and
\(F(v)+F(u)\) within \(3\eta\).  Distinct Weyl axes have phase-optimized
distance \(\sqrt2\), so \(\eta<1/3\) makes \(F\) additive.  Comparing
the exact commutator phases of \(A_v,A_u\) with those of their nearby
Weyls costs at most \(4\eta\).  The least nontrivial \(p\)-th-root chord
is \(2\sin(\pi/p)\), so the stated bound makes \(F\) trace-symplectic.

Choose an additive Clifford \(K_0\) inducing \(F\).  Then
\(K_0^\dagger L\) approximately fixes every Weyl axis, so all its Weyl
commutators are within \(\eta\) of scalars.  The Fourier observation puts
it within \(\eta\) of a Weyl; multiplying that Weyl into \(K_0\) gives
the required Clifford.
\end{proof}

Two further estimates turn local Clifford frames into the stated global
decomposition.  The first separates distinct Clifford images of a stabilizer
state by a uniform gap.

\begin{lemma}[Quantized stabilizer overlap]
\label{lem:stabilizer-overlap-gap}
Let \(\lvert\psi\rangle,\lvert\phi\rangle\) be stabilizer states on \(n\)
qudits of local dimension \(q=p^e\), with label groups
\(L_\psi,L_\phi\) and stabilizer characters \(\chi_\psi,\chi_\phi\).  Then
\[
 \lvert\langle\psi\vert\phi\rangle\rvert^2=
 \begin{cases}
  \lvert L_\psi\cap L_\phi\rvert/q^n,
   &\text{if \(\chi_\psi=\chi_\phi\) on \(L_\psi\cap L_\phi\)},\\
  0,&\text{otherwise.}
 \end{cases}
\]
In particular, a product Clifford \(K\) either preserves the ray of
\(\psi\), or
\[
 \varepsilon(K)\geq d_p:=\bigl(2-2p^{-1/2}\bigr)^{1/2}.
\]
\end{lemma}

\begin{proof}
Weyl orthogonality and the stabilizer expansions give
\[
 \lvert\langle\psi\vert\phi\rangle\rvert^2
 =q^{-n}\sum_{\ell\in L_\psi\cap L_\phi}
   (\chi_\psi\overline{\chi_\phi})(\ell).
\]
The summand is a character of the intersection, so the sum is its order when
the character is trivial and zero otherwise.  Since the intersection has
\(p\)-power order, every nonunit overlap has modulus at most \(p^{-1/2}\).
Apply this to \(\phi=K\psi\), which is again a stabilizer state.
\end{proof}

The second estimate controls all residual generators at once.  Its proof is
included because it is a load-bearing part of the headline quantitative
theorem; Appendix~\ref{sec:two-uniform} gives stronger variants under other
hypotheses.

\begin{proposition}[Residual stability]
\label{prop:main-residual-stability}
Let \(\lvert\psi\rangle\) be a stabilizer \(\AME(2m,q)\) state with
\(m\geq2\), and let \(V=\bigotimes_je^{ih_j}\), where every \(h_j\) is
traceless Hermitian of spectral spread at most \(s\leq\pi\).  Put
\[
 c(s)=\max\left\{1-\frac{s^2}{12},\frac4{\pi^2}\right\}.
\]
If \(D^2=\sum_j\lVert h_j\rVert_F^2\leq q/4\), then
\[
 D\leq2\sqrt{q/c(s)}\,\varepsilon(V)
 \leq\pi\sqrt q\,\varepsilon(V).
\]
\end{proposition}

\begin{proof}
Expand each factor in the normalized Weyl basis as
\(e^{ih_j}=\sum_va_j(v)W_v\), and put
\(\eta_j=1-|a_j(0)|^2\) and \(H=\sum_j\eta_j\).  Choose an \(m\)-party set
\(A\).  The AME support formula makes restriction of the stabilizer label
group \(L\) to either \(A\) or \(A^c\) a bijection.  Hence the vectors
\[
 u(\ell)=\prod_{j\in A}|a_j(\ell_j)|,
 \qquad
 v(\ell)=\prod_{j\notin A}|a_j(\ell_j)|
 \quad(\ell\in L)
\]
are unit vectors.  Expanding the stabilizer expectation and taking moduli
termwise gives \(\varepsilon(V)^2\geq\lVert u-v\rVert^2\).

For each \(j\notin A\) and Weyl label \(w\), let \(\ell^{j,w}\) be the
unique stabilizer label supported on \(A\cup\{j\}\) whose \(j\)-coordinate
is \(w\).  The support formula makes this parametrization bijective and gives
\[
 \sum_{w\ne0}v(\ell^{j,w})^2
 \geq\eta_j(1-H),
 \qquad
 \sum_{w\ne0}u(\ell^{j,w})^2
 \leq\prod_{i\in A}\eta_i.
\]
The corresponding families are disjoint as \(j\) varies.  Applying
\((x-y)^2\geq x^2/2-y^2\), and then repeating the estimate with \(A\) and
\(A^c\) exchanged, yields
\[
 \lVert u-v\rVert^2
 \geq\bigl(\tfrac12(1-H)-\eta^{m-1}\bigr)H,
 \qquad \eta=\max_j\eta_j.
\]
For \(m\geq3\) and \(H\leq1/4\), the coefficient is at least \(5/16\).
For \(m=2\), retaining the four pair products before the maximum estimate
gives \(H/2-3P\), where \(P\leq H^2/4\), again at least \(5H/16\).
Thus in all cases \(\varepsilon(V)^2\geq H/4\).

It remains to compare \(H\) with \(D\).  If \(X\) is uniform on the
eigenvalues of \(h_j\), then
\[
 \eta_j=1-|\mathbb E e^{iX}|^2
 =\mathbb E[1-\cos(X-X')]
 \geq c(s)\frac{\lVert h_j\rVert_F^2}{q},
\]
using both \(\cos x\leq1-x^2/2+x^4/24\) for \(|x|\leq s\) and
\(1-\cos x\geq2x^2/\pi^2\) for \(|x|\leq\pi\).  Conversely
\(\eta_j\leq\lVert h_j\rVert_F^2/q\), so \(H\leq D^2/q\leq1/4\).
Combining the inequalities gives
\(\varepsilon(V)^2\geq c(s)D^2/(4q)\), which proves the claim.
\end{proof}

At \(\varepsilon=0\), the first two lemmas in this section recover the
three-cleanable-region Clifford conclusion of
\cite[Lemma~5]{PastawskiYoshida2015}.  Their quantitative combination, and
the constants in (5.1), are what will be used here.

\begin{theorem}[Cleaning-based global rounding]
\label{thm:cleaning-global-rounding}
Let \(\lvert\psi\rangle\) be a stabilizer \(\AME(2m,q)\) state,
\(q=p^e\), \(m\geq2\), and \(n=2m\).  Put
\[
 d_p=(2-2p^{-1/2})^{1/2},
\]
and
\[
 R^{\mathrm{clean}}_{n,q}=\min\left\{
  \frac{\tau_p}{8},\frac1{4\sqrt{2q}},
  \frac1{8\pi\sqrt n},
  \frac{d_p}{1+4\pi\sqrt n}
 \right\}.                                           \tag{5.3}
\]
If \(\varepsilon(U)<R^{\mathrm{clean}}_{n,q}\), then, up to a global
phase,
\[
 U=g\bigotimes_{i=1}^n e^{ih_i},\qquad g\in G(\psi),
\]
where the \(h_i\) are traceless Hermitian, have spectral spread at most
\(\pi\), and satisfy
\[
 \left(\sum_i\lVert h_i\rVert_F^2\right)^{1/2}
 \leq\pi\sqrt q\,\varepsilon(U).                     \tag{5.4}
\]
Before the exact-symmetry step, the recovered local Cliffords obey
\[
 q^{-1/2}\min_{|z|=1}\lVert U_i-zK_i\rVert_{\mathrm{HS}}
 \leq8\varepsilon(U)                                 \tag{5.5}
\]
at every party.  For fixed \(q\), (5.3) has order \(n^{-1/2}\),
conditional on existence of the stated AME states.
\end{theorem}

\begin{proof}
\medskip\noindent\emph{Local Clifford frames.}
Choose party \(i\) as the logical leg and write the AME Choi form as
\(q^{-1/2}\sum_a|a\rangle\otimes V_i|a\rangle\).  The other local
factors form a transversal physical unitary \(T_i\), and vectorization
gives a logical unitary \(L_i\), obtained from \(U_i\) by transpose and
inverse, with \(E(T_i,L_i)=\varepsilon(U)\).  Lemmas
\ref{lem:quantitative-cleaning-commutator} and
\ref{lem:nested-weyl-rounding}, with \(\eta=8\varepsilon(U)\), give
(5.5); transpose and inverse preserve the additive Clifford group.

\medskip\noindent\emph{Centered residual generators.}
Put \(K=\bigotimes_iK_i\).  For each \(i\), choose a phase \(z_i\)
attaining the minimum in (5.5), and let \(a_i\) be the principal Hermitian
logarithm defined by
\[
 z_i^{-1}K_i^\dagger U_i=e^{ia_i}.
\]
Set \(\bar a_i=q^{-1}\operatorname{Tr}(a_i)\) and \(h_i=a_i-\bar a_iI\).
Rephase \(K_i\) by \(z_ie^{i\bar a_i}\); this changes
\(K=\bigotimes_iK_i\) only by a global phase and does not change its
projective Clifford action.  After this rephasing,
\(K_i^\dagger U_i=e^{ih_i}\), each \(h_i\) is traceless, centering does not
increase its Frobenius norm, and its spectral spread is unchanged.  Let
\(D^2=\sum_i\lVert h_i\rVert_F^2\).  The chord--angle bound gives
\[
 D^2\leq\frac{\pi^2q}{4}\sum_i
  \left(q^{-1/2}\min_{|z|=1}
   \lVert U_i-zK_i\rVert_{\mathrm{HS}}\right)^2
 \leq16\pi^2qn\varepsilon(U)^2.                      \tag{5.6}
\]
The second clause of (5.3) puts every spectrum in an arc of length
\(\pi\); the third gives \(D\leq\sqrt q/2\).

\medskip\noindent\emph{Exact branch selection.}
Two-uniformity gives
\(\lVert(K^\dagger U-I)\psi\rVert\leq D/\sqrt q
\leq4\pi\sqrt n\,\varepsilon(U)\).  The fourth clause and the triangle
inequality put \(K\psi\) less than \(d_p\) from \(\psi\), so
Lemma~\ref{lem:stabilizer-overlap-gap} makes \(K\) an exact symmetry.
Apply Proposition~\ref{prop:main-residual-stability} to obtain (5.4).
\end{proof}

The theorem already contains a genuinely collective estimate once the
exact branch has been selected.  Writing \(g=\bigotimes_i g_i\) in the
factorization above,
\[
 \sum_i\left(q^{-1/2}\min_{|z|=1}
   \lVert U_i-zg_i\rVert_{\mathrm{HS}}\right)^2
 \leq \frac{D^2}{q}
 \leq \pi^2\varepsilon(U)^2.                         \tag{5.7}
\]
The first inequality is the elementary chord bound
\(\lVert e^{ih_i}-I\rVert_{\mathrm{HS}}\leq\lVert h_i\rVert_F\).
Thus the \(\sqrt n\) loss is absent after exact-symmetry recovery.

The same estimate applies to intertwiners between two states as soon as one
exact intertwiner is fixed.  This separates the cost of finding the correct
orbit from stability within that orbit.

\begin{corollary}[Relative intertwiner rounding]
\label{cor:relative-intertwiner-rounding}
Let \(\lvert\psi\rangle,\lvert\phi\rangle\) be stabilizer
\(\AME(2m,q)\) states as in
Theorem~\ref{thm:cleaning-global-rounding}, and suppose that
\(B=\bigotimes_iB_i\) is an exact product-unitary intertwiner from
\(\psi\) to \(\phi\).  For \(U=\bigotimes_iU_i\), put
\[
 \varepsilon_{\psi\to\phi}(U)=
 \min_{|z|=1}\lVert U\psi-z\phi\rVert.
\]
If \(\varepsilon_{\psi\to\phi}(U)<\tau_p/8\), then each \(U_i\) is within
normalized Hilbert--Schmidt distance
\(8\varepsilon_{\psi\to\phi}(U)\) of \(B_iK_i\) for a one-qudit Clifford
\(K_i\).  If, in addition,
\(\varepsilon_{\psi\to\phi}(U)<R^{\mathrm{clean}}_{n,q}\), then, up to
phase,
\[
 U=H\bigotimes_i e^{ih_i},
 \qquad H\psi\in\mathbb C\phi,
 \qquad
 \left(\sum_i\lVert h_i\rVert_F^2\right)^{1/2}
 \leq\pi\sqrt q\,\varepsilon_{\psi\to\phi}(U),
\]
where \(H\) is an exact product-unitary intertwiner and the \(h_i\) are
traceless Hermitian of spectral spread at most \(\pi\).  After branch
selection,
\[
 \sum_i\left(q^{-1/2}\min_{|z|=1}
   \lVert U_i-zH_i\rVert_{\mathrm{HS}}\right)^2
 \leq\pi^2\varepsilon_{\psi\to\phi}(U)^2.
\]
Thus the stability constants are independent of the party count; the
certified radius remains \(R^{\mathrm{clean}}_{n,q}\).
\end{corollary}

\begin{proof}
Set \(W=B^\dagger U=\bigotimes_iB_i^\dagger U_i\).  Since \(B\psi\) and
\(\phi\) span the same ray,
\(\varepsilon(W)=\varepsilon_{\psi\to\phi}(U)\).  The first paragraph of
the proof of Theorem~\ref{thm:cleaning-global-rounding} gives the local
\(8\varepsilon\) estimate under \(\varepsilon(W)<\tau_p/8\).  Under the
stronger radius hypothesis, apply the theorem and (5.7) to \(W\).  If
\(W=g\bigotimes_i e^{ih_i}\) with \(g\in G(\psi)\), then \(H=Bg\) is an
exact intertwiner from \(\psi\) to \(\phi\), and left multiplication by
each \(B_i\) preserves normalized Hilbert--Schmidt distance.  This proves
all three estimates.  Exact rigidity also makes every \(B_i\) Clifford, but
that fact is not needed for the reduction.
\end{proof}

This does not enlarge the entry radius in (5.3).  Before the branch is
known, cleaning supplies only the separate bounds (5.5) for locally
chosen Cliffords \(K_i\).  The \(n\) encoded implementation errors are
unitary reshapings of the same global defect vector, not projections onto
orthogonal sectors, so their squared bounds cannot be summed by a Bessel
argument.  A radius-improving collective estimate would itself imply
constant-radius branch selection: the chord--angle inequality would make
\(\sum_i\lVert h_i\rVert_F^2=O(q\varepsilon^2)\), and the uniform
stabilizer gap would force \(\bigotimes_iK_i\) to be exact.  The missing
input is therefore compatibility of the recovered discrete Clifford data,
not a sharper final norm summation.

The linear part of that compatibility can in fact be recovered without a
party-count loss.

\begin{proposition}[Robust compatibility of the transition maps]
\label{prop:robust-linear-atlas}
Let \(\lvert\psi\rangle\) be a stabilizer \(\AME(2m,q)\) state,
\(q=p^e\), and let \(U=\bigotimes_iU_i\) be a product unitary.  Suppose
\[
 \varepsilon(U)<
 R_q^{\mathrm{trans}}
 :=\min\left\{\frac{\tau_p}{8},
   \frac{2\sin(\pi/p)}{4+64\sqrt q}\right\}.
\]
The cleaning and Fourier lemmas give local Cliffords \(K_i\) satisfying
(5.5).  Let \(F_i\) be their symplectic maps.  Then
\[
 F_jM_A(i,j)=M_A(i,j)F_i
\]
for every \((m+1)\)-set \(A\) and \(i,j\in A\).  Consequently
\(F=\bigoplus_iF_i\) preserves the stabilizer label group \(L\), and some
product-Pauli correction \(W_vK\) is an exact symmetry of \(\psi\).
\end{proposition}

\begin{proof}
Fix \(A\), distinct \(i,j\in A\), and put
\(B=A^c\cup\{i,j\}\).  Both \(A\) and \(B\) have size \(m+1\), while
\(A\cap B=\{i,j\}\).  Choose \(\ell\in L(A)\) and \(r\in L(B)\).
Using the stabilizer lifts \(\widetilde W_\ell\) and
\(\widetilde W_r\) from \((2.2)\), set
\[
 C=[U\widetilde W_\ell U^\dagger,\widetilde W_r]
  =[UW_\ell U^\dagger,W_r],
\]
where the second equality holds because scalar lift phases cancel from a
commutator.
After choosing the phase of \(U\) realizing its defect,
\(\lVert(U\widetilde W_\ell U^\dagger-I)\psi\rVert
\leq2\varepsilon(U)\), and the same holds with the adjoint.  Since
\(\widetilde W_r\psi=\psi\),
\[
 \lVert(C-I)\psi\rVert\leq4\varepsilon(U).             \tag{5.8}
\]

The commutator is supported only on \(A\cap B=\{i,j\}\).
Choose phases so that
\(\lVert U_s-K_s\rVert_{\mathrm{HS}}\leq8\sqrt q\,\varepsilon(U)\).
Replacing \(U_sW_{\ell_s}U_s^\dagger\) by its \(K_s\)-conjugate costs at
most \(16\sqrt q\,\varepsilon(U)\) in operator norm, and replacing its
commutator with \(W_{r_s}\) costs at most twice that.  On the two sites,
the total cost is therefore at most
\(64\sqrt q\,\varepsilon(U)\).  The Clifford commutator is a scalar
\(p\)-th root \(\zeta\), so (5.8) gives
\[
 |\zeta-1|\leq(4+64\sqrt q)\varepsilon(U)<2\sin(\pi/p).
\]
Hence \(\zeta=1\).

Write \(a=\ell_i\), \(b=r_i\),
\(\ell_j=M_A(i,j)a\), and \(r_j=M_B(i,j)b\).
The identity \(\zeta=1\), for all \(a,b\), says
\[
 [F_i a,b]+[F_jM_A(i,j)a,M_B(i,j)b]=0.                \tag{5.9}
\]
Here brackets denote the trace-symplectic pairing.  Isotropy of \(L\),
applied to the element of \(L(A)\) whose \(i\)-coordinate is \(F_i a\),
gives
\[
 [F_i a,b]+[M_A(i,j)F_i a,M_B(i,j)b]=0.
\]
Subtracting this from (5.9), using surjectivity of \(M_B(i,j)\) and
nondegeneracy of the pairing, proves the transition equation.
The exact transition-map theorem gives \(FL=L\), and
Lemma~\ref{lem:pauli-phase-correction} supplies \(W_v\).
\end{proof}

The uncontrolled Pauli correction does not obstruct the logical action of a
chosen encoder.

\begin{corollary}[Logical Clifford rounding]
\label{cor:logical-clifford-rounding}
Under the hypotheses and radius of
Proposition~\ref{prop:robust-linear-atlas}, choose party \(i\) as the logical
leg, and let \(L_i\) be the logical unitary induced by the other \(2m-1\)
factors as in the proof of Theorem~\ref{thm:cleaning-global-rounding}.  There
is an exactly transversally realizable logical Clifford \(L_i'\) such that
\[
 q^{-1/2}\min_{|z|=1}\lVert L_i-zL_i'\rVert_{\mathrm{HS}}
 \leq8\varepsilon(U).
\]
\end{corollary}

\begin{proof}
Write the exact symmetry supplied by the proposition as
\(W_vK\), where \(K=\bigotimes_jK_j\).  The projection of the stabilizer
label group onto the Weyl plane at party \(i\) is surjective: this already
holds for every minimum-support subgroup containing \(i\).  Multiply \(W_vK\)
by a stabilizer whose \(i\)-th label is \(-v_i\).  The result is still an exact
product-Clifford symmetry, and its factor at \(i\) is projectively \(K_i\).
The Choi correspondence therefore makes its other factors an exact
transversal implementation of
\(L_i'=(K_i^{\mathsf T})^{-1}\).  The logical unitary \(L_i\) is the same
transpose--inverse transform of \(U_i\), up to phase, so (5.5) proves the
displayed bound.
\end{proof}

The Pauli correction in the proposition is essential.  All localized
commutators are unchanged when a recovered \(K_i\) is multiplied by a
local Weyl: that multiplication changes conjugation phases, and the
commutator cancels them.  Equivalently, the transition maps see \(L\)
but not its stabilizer character.  The correction \(v\) can therefore be
nonzero modulo \(L\), in which case \(W_vK\) need not be locally close to
\(U\).  Exact affine compatibility cannot be obtained from these
commutator tests, and Proposition~\ref{prop:robust-linear-atlas} alone
does not improve (5.3).  Controlling the character discrepancy is the
same remaining branch-selection problem isolated by the collective
estimate above.

The exact formula, rather than a fixed-dimension slogan, controls scaling.
Uniformly over prime powers \(q=p^e\) and existing tensors,
\[
 R^{\mathrm{clean}}_{n,q}
 =\Theta\!\left(\min\{p^{-1},q^{-1/2},n^{-1/2}\}\right),              \tag{5.10}
\]
with absolute implied constants.  Indeed
\(\sin(\pi/p)=\Theta(p^{-1})\), while \(d_p\) is bounded above and below
by positive absolute constants.  An asymptotic use of (5.10) must specify
an existing family.

For the explicit generalized and extended generalized Reed--Solomon
construction \cite{SeroussiRoth1986}, together with the standard
MDS-to-AME construction \cite{RaissiGogolinRieraAcin2018},
\(n=2m\leq q+1\), so the party-count term cannot be smaller
than the local-dimension term.  Hence
\[
 R^{\mathrm{clean}}_{n,q}
 =\begin{cases}
   \Theta(q^{-1}),&q=p,\\
   \Theta(q^{-1/2}),&q=p^e,\ e\geq2.
  \end{cases}                                                        \tag{5.11}
\]
On a length-saturating subfamily \(n\asymp q\), these are
\(\Theta(n^{-1})\) and \(\Theta(n^{-1/2})\), respectively.  For fixed
characteristic and growing extension degree, the second line applies.
The earlier fixed-\(q\), \(n^{-1/2}\) reading of the exact minimum remains
a conditional formal slice only; it is not an AME-existence claim.

If the sharper local constant and the old global generator budget are
needed, the same proof gives the conclusion of
Corollary~\ref{cor:approximate-decomposition} under
\[
 \varepsilon(U)<\min\left\{
  \frac{\tau_p}{8},\frac1{16\pi n\sqrt q}
 \right\}.                                           \tag{5.12}
\]
Indeed (5.5), the operator-norm bound by Hilbert--Schmidt norm, and the
principal logarithm give
\(\sum_i\lVert h_i\rVert_{\mathrm{op}}
\leq8\pi n\sqrt q\,\varepsilon(U)<1/2\); the sharp local theorem then
gives \(D\leq\sqrt{6q/5}\,\varepsilon(U)\).

The cleaning route supersedes both explicit radii above on every admissible
parameter pair where its displayed minimum is larger, and it has no
boundary at \(q=4\).  The earlier theorems remain useful because they expose the exact
cost of reading one marginal and of aggregating orthogonal minimum-support
sectors.  The new \(n^{-1/2}\) scale is a proved lower bound, not an
optimality claim; no product-unitary family is known to saturate the
collective frame-error estimate (5.6).
